\lecture{8.10}{Time-to-Event and Competing Risks}
\label{lec:8-10}

\subsection{Longitudinal Data With Event Times}

Time-to-event data can be represented using the same longitudinal structure as
time-varying treatments. Let $T$ be the event time of interest and let $C$ be a
right-censoring time. The observed time is
\[
  \widetilde T=\min(T,C),
\]
with event indicator
\[
  \Delta=\ind(T\le C).
\]
For causal questions, one often wants the counterfactual cumulative incidence
\[
  \Prob\{T(a)\le t\}
\]
or its longitudinal analogue under a treatment or censoring intervention.

Discretize time into intervals $1,\ldots,t$. Define
\[
  L_j=\ind(\text{event has occurred by time } j),
\]
and let $A_j$ indicate remaining uncensored through the relevant interval. Then
\[
  \Prob\{T(a)\le t\}=\Prob\{L_t(a)=1\}.
\]
The event process $L_j$ is a counting process with an absorbing state: once
$L_j=1$, all later $L_s$ for $s>j$ must also equal 1.

\subsection{Censoring}

Right censoring can be viewed as a time-varying process that prevents later
outcomes from being observed. If $A_j=1$ denotes not being censored through time
$j$, then estimating a full-data event probability can be framed as estimating
the outcome under the intervention
\[
  A_1=A_2=\cdots=A_t=1,
\]
that is, the intervention that eliminates censoring.

\begin{assumptionbox}
Identification under censoring requires sequential exchangeability for the
censoring process:
\[
  L_t(\bar a) \ci A_j \given \Lbar_j,\Abar_{j-1}
  \qquad j=1,\ldots,t.
\]
In words, among people with the same observed history, remaining uncensored is
independent of the counterfactual event status that would be observed under the
no-censoring intervention.
\end{assumptionbox}

\subsection{Sequential Regression for Survival Outcomes}

Consider estimating $\Prob\{L_3(1,1,1)=1\}$, the cumulative incidence by time 3
under no censoring. A sequential regression algorithm works backward:

\begin{enumerate}
  \item Estimate
  \[
    \Qbar_3(\Lbar_2)=\E(L_3\given \Abar_3=\mathbf{1},\Lbar_2).
  \]
  If $L_2=1$, then $\Qbar_3(\Lbar_2)=1$ by the absorbing event structure.
  Therefore only histories with $L_2=0$ require a regression model.
  \item Build a pseudo-outcome $\Qbar_{3,n}$ that equals 1 for those with
  $L_2=1$ and equals the fitted probability for those with $L_2=0$.
  \item Estimate
  \[
    \Qbar_2(\Lbar_1)=\E\{\Qbar_3(\Lbar_2)\given \Abar_2=\mathbf{1},\Lbar_1\}.
  \]
  If $L_1=1$, then $\Qbar_2(\Lbar_1)=1$.
  \item Continue backward and average the final prediction over baseline
  covariates.
\end{enumerate}

\begin{lstlisting}[language=R]
# L1, L2, L3 are indicators of event by times 1, 2, 3.
# A1, A2, A3 are indicators of remaining uncensored.

full_data$Q3n <- NA
full_data$Q3n[full_data$L2 == 1] <- 1

fit_data_Q3 <- subset(
  full_data,
  A1 == 1 & A2 == 1 & A3 == 1 & L2 == 0
)
fit_Q3 <- glm(L3 ~ L0, data = fit_data_Q3, family = binomial())

pred_data <- full_data[full_data$L2 == 0, ]
full_data$Q3n[full_data$L2 == 0] <-
  predict(fit_Q3, newdata = pred_data, type = "response")

full_data$Q2n <- NA
full_data$Q2n[full_data$L1 == 1] <- 1

fit_data_Q2 <- subset(full_data, A1 == 1 & A2 == 1 & L1 == 0)
fit_Q2 <- glm(Q3n ~ L0, data = fit_data_Q2, family = binomial())

pred_data <- full_data[full_data$L1 == 0, ]
full_data$Q2n[full_data$L1 == 0] <-
  predict(fit_Q2, newdata = pred_data, type = "response")

fit_data_Q1 <- subset(full_data, A1 == 1)
fit_Q1 <- glm(Q2n ~ L0, data = fit_data_Q1, family = binomial())
full_data$Q1n <- predict(fit_Q1, newdata = full_data, type = "response")

cuminc_3 <- mean(full_data$Q1n)
\end{lstlisting}

\subsection{Competing Risks}

In competing risks, the first event may be one of several types. Let
$J\in\{1,\ldots,m\}$ denote event type. The observed event indicator may take
values
\[
  \Delta=0 \text{ for censoring},\qquad
  \Delta=j \text{ for an event of type } j.
\]
For discrete time, define
\[
  L_{kj}=\ind(\text{event of type } j \text{ has occurred by time } k),
\]
and collect these into
\[
  L_k=(L_{k1},\ldots,L_{km}).
\]
The type-$j$ cumulative incidence is
\[
  \Prob\{T(a)\le t,J(a)=j\}=\Prob\{L_{tj}(a)=1\}.
\]

The absorbing structure is more restrictive than in the single-event case. If a
type 1 event occurred by time $k$, then type 1 remains occurred at all later
times and no other first-event type can occur later. For two event types:
\[
  L_{k1}=1 \Rightarrow L_{s1}=1,\ L_{s2}=0
  \quad\text{for }s>k.
\]

\subsection{Counterfactual Cumulative Incidence}

The counterfactual cumulative incidence under treatment or no censoring is a
standard longitudinal counterfactual mean. The outcome is the event indicator by
time $t$, and the intervention may set treatment, censoring, or both.

For example, with no censoring through time $t$,
\[
  \psi_j(t)=\E\{L_{tj}(\mathbf{1})\}
  =
  \Prob\{T(\mathbf{1})\le t,J(\mathbf{1})=j\}.
\]
Contrasts of $\psi_j(t)$ across treatment regimes are causal contrasts of
cumulative incidence curves.

\subsection{Practical Estimation Issues}

Survival and competing-risk estimation require the analyst to respect deterministic
history restrictions:
\begin{itemize}
  \item do not fit event regressions among people who already had the event when
  the outcome is already known;
  \item set the pseudo-outcome to 1 after the target event has occurred;
  \item set the pseudo-outcome to 0 after a competing first event has occurred
  when estimating the cumulative incidence of a different event type;
  \item define censoring indicators carefully so that future outcomes are not
  used after censoring.
\end{itemize}

\begin{keyidea}
Time-to-event data become longitudinal data after discretization. The main extra
work is enforcing the absorbing-state logic of event histories and, for competing
risks, the mutually exclusive nature of first-event types.
\end{keyidea}
